\section{Full Mermaid Compatibility}
\label{sec:mermaid-compat}

The compiler targets \textbf{complete compatibility with all 22 Mermaid diagram types}.
Every valid Mermaid file must parse, validate, and render---producing output that is
semantically equivalent to Mermaid's but visually superior.

\subsection{The 22 Diagram Types}

Mermaid supports 22 diagram types as of 2026. We organize them into five families
(Section~\ref{sec:family-taxonomy}) and assign coverage tiers:

\begin{table}[h]
\centering\small
\caption{Mermaid diagram types: families and coverage tiers}
\begin{tabular}{llll}
\toprule
\textbf{Mermaid Type} & \textbf{Family} & \textbf{Tier} & \textbf{Status} \\
\midrule
flowchart       & Node-link / graph     & 0 & Built (Flow grammar) \\
sequence        & UML / software        & 0 & Built (Sequence grammar) \\
gantt           & Timeline / project    & 0 & Built (Timeline grammar) \\
timeline        & Timeline / project    & 0 & Built (Timeline grammar) \\
mindmap         & Tree / hierarchy      & 0 & Built (Tree grammar) \\
\midrule
class           & UML / software        & 1 & Planned \\
state           & UML / software        & 1 & Planned \\
erDiagram       & UML / software        & 1 & Planned \\
C4Context       & Node-link / graph     & 1 & Planned \\
C4Container     & Node-link / graph     & 1 & Planned \\
C4Component     & Node-link / graph     & 1 & Planned \\
C4Dynamic       & Node-link / graph     & 1 & Planned \\
\midrule
pie             & Charts                & 2 & Planned \\
xychart-beta    & Charts                & 2 & Planned \\
quadrantChart   & Charts                & 2 & Planned \\
radar           & Charts                & 2 & Planned \\
\midrule
sankey-beta     & Node-link / graph     & 3 & Planned \\
requirement     & Node-link / graph     & 3 & Planned \\
gitGraph        & Node-link / graph     & 3 & Planned \\
block-beta      & Node-link / graph     & 3 & Planned \\
architecture    & Node-link / graph     & 3 & Planned \\
treemap         & Tree / hierarchy      & 3 & Planned \\
kanban          & Timeline / project    & 3 & Planned \\
journey         & Timeline / project    & 3 & Planned \\
packet          & Node-link / graph     & 3 & Planned \\
\bottomrule
\end{tabular}
\label{tab:mermaid-coverage}
\end{table}

\subsection{Compatibility Strategy}

\subsubsection{Syntax Fidelity}

Each Mermaid grammar is parsed faithfully. Our parsers accept the exact syntax documented at
\texttt{mermaid.js.org}~\cite{mermaiddocs2024}, including:

\begin{itemize}
  \item Indentation-sensitive formats (timeline, mindmap, gantt)
  \item Arrow syntax variants (\texttt{-->}, \texttt{-.->}, \texttt{==>}, etc.\ in flowchart)
  \item Section/participant declarations
  \item Frontmatter (\texttt{---}$\ldots$\texttt{---}) and init directives (\texttt{\%\%\{init\}\%\%})
  \item All node shape syntaxes (round, diamond, hexagon, etc.)
\end{itemize}

\subsubsection{Semantic Mapping}

Mermaid's syntax maps to our Domain IRs via grammar-specific lowering rules:

\begin{description}
  \item[flowchart $\to$ Flow IR] Nodes become Flow nodes; edges become Flow edges;
        subgraphs become lanes/clusters.
  \item[sequence $\to$ Sequence IR] Participants, messages, activations, and fragments
        map directly.
  \item[gantt/timeline $\to$ Timeline IR] Sections become tracks; tasks become activities;
        milestones map directly.
  \item[mindmap $\to$ Tree IR] Indentation hierarchy becomes parent-child tree structure.
\end{description}

\subsubsection{Directive and Frontmatter Hooks}

Mermaid's configuration mechanisms become our extension points:

\begin{lstlisting}[language=yaml,caption={Mermaid frontmatter as theme/extension hook}]
---
title: Architecture Overview
theme: professional-dark
config:
  flowchart:
    nodeSpacing: 60
    rankSpacing: 80
---
flowchart LR
    A[Client] --> B[API Gateway]
    B --> C[Service Mesh]
\end{lstlisting}

The \texttt{theme} field selects our theme system. The \texttt{config} block maps to
grammar-specific theme extensions. Init directives (\texttt{\%\%\{init: \{...\}\}\%\%})
provide inline configuration.

\subsection{Reuse of Existing Kernels}

Most new diagram types reuse existing layout kernels:

\begin{description}
  \item[Node-link kernel] (Sugiyama/layered layout) serves: flowchart, C4 variants,
        architecture, block, requirement, gitGraph, sankey.
  \item[Tree kernel] (Buchheim--J\"unger--Leipert) serves: mindmap, treemap.
  \item[Timeline kernel] (track-based layout) serves: gantt, timeline, journey, kanban.
  \item[Sequence kernel] (order-deterministic columns) serves: sequence diagrams.
  \item[Chart kernel] (NEW grammar-of-graphics layer) serves: pie, xychart, quadrant, radar.
\end{description}

The chart family is the only genuinely new kernel required (Section~\ref{sec:chart-family}).
All other new types are primarily new \emph{parsers} wiring Mermaid syntax to existing
Domain IR shapes and layout engines.
